% Source fingerprint: 326ba469d966e476facdf6f0d0b0b875b3f04f1cca131e889bf24c095f997433
% T11: raw TeX cells preserved; no numeric highlighting.
\begin{table}[htbp]
\centering\small\renewcommand{\arraystretch}{1.18}\setlength{\tabcolsep}{4pt}
\caption[零边界、零延拓与零迹]{零边界、零延拓与零迹。这是不同边界条件的数学比较。各行只使用本文已建立的方向，不把零延拓或逐点边界值无条件等同于零迹。}
\label{b1:tab:visual-t11}
\begin{tabularx}{\linewidth}{@{}>{\raggedright\arraybackslash}p{3.1cm}>{\raggedright\arraybackslash}p{4.7cm}>{\raggedright\arraybackslash}X@{}}
\toprule
条件 & 可以得到的结论 & 适用边界 \\
\midrule
\(u\in W_0^{1,p}(\Omega)\) & 按 \(C_c^\infty\) 的 \(W^{1,p}\) 闭包定义；零延拓属于全空间 Sobolev & 本章对 \(1\le p\le\infty\) 使用此定义；零延拓的反向结论不在任意域无条件使用 \\
连续代表在边界为零 & 有界开集、有限 \(p\) 时推出 \(u\in W_0^{1,p}\) & 不要求该充分条件中的边界光滑；但一般 Sobolev 元素未必有连续代表 \\
区间端点值为零 & 有限 \(p\) 时等价于 \(W_0^{1,p}(a,b)\) & 使用绝对连续代表；无穷端点的闭包条件更强 \\
迹算子取零 & 有界 Lipschitz 区域、有限 \(p\) 时，本章已证 \(\ker\operatorname{Tr}=W_0^{1,p}\) & 这里的迹由内部逼近构造，不是任意可测代表的边界赋值 \\
\(p=\infty\) 的连续零边界 & 不足以推出 \(W_0^{1,\infty}\) & 例如区间上的二次函数，导数不能被内部紧支集光滑导数一致逼近 \\
\bottomrule
\end{tabularx}
\end{table}
